\Article{33}
Wegen der Bedeutung des Gegenstands und damit alles für die Rechnung bereitliegt, schien es angebracht, hier auch die expliziten Formeln zur Bestimmung der Koeffizienten \(A',A'',A'''\) etc., \(B'',B'''\) etc. etc. anzugeben. Diese Rechnung lässt sich auf zwei Arten ausführen, da sowohl beim Einsetzen der aus III entnommenen Werte von \(u^0,u',u''\) etc. in II als auch beim Einsetzen der Werte von \(\xi,\eta,\zeta\) etc. aus II in III Identitäten entstehen müssen. Die erste Art liefert die folgenden Formelsysteme:
\[
\begin{aligned}
\frac{\mathfrak{B}^0}{\mathfrak{A}^0}+A'&=0\\[4pt]
\frac{\mathfrak{C}^0}{\mathfrak{A}^0}+\frac{\mathfrak{C}'}{\mathfrak{B}'}.A'+A''&=0\\[4pt]
\frac{\mathfrak{D}^0}{\mathfrak{A}^0}+\frac{\mathfrak{D}'}{\mathfrak{B}'}.A'+\frac{\mathfrak{D}''}{\mathfrak{C}''}.A''+A'''&=0
\end{aligned}
\]
etc., woraus sich \(A',A'',A'''\) etc. ergeben.
\[
\begin{aligned}
\frac{\mathfrak{C}'}{\mathfrak{B}'}+B''&=0\\[4pt]
\frac{\mathfrak{D}'}{\mathfrak{B}'}+\frac{\mathfrak{D}''}{\mathfrak{C}''}.B''+B'''&=0
\end{aligned}
\]
etc., woraus sich \(B'',B'''\) etc. ergeben.
\[
\frac{\mathfrak{D}''}{\mathfrak{C}''}+C'''=0
\]
etc., woraus sich \(C'''\) etc. ergeben. Und so fort.

Die andere Art liefert diese Formeln:
\[
\mathfrak{A}^0A'+\mathfrak{B}^0=0
\]
woraus sich \(A'\) ergibt.
\[
\begin{array}{rllll}
\mathfrak{A}^0A''&+\mathfrak{B}^0B''&+\mathfrak{C}^0&=0\\
&\mathfrak{B}'B''&+\mathfrak{C}'&=0
\end{array}
\]
woraus sich \(B''\) und \(A''\) ergeben.
\SourceLeaf{046}
\[
\begin{array}{rlllll}
\mathfrak{A}^0A'''&+\mathfrak{B}^0B'''&+\mathfrak{C}^0C'''&+\mathfrak{D}^0&=0\\
&\mathfrak{B}'B'''&+\mathfrak{C}'C'''&+\mathfrak{D}'&=0\\
&&\mathfrak{C}''C'''&+\mathfrak{D}''&=0
\end{array}
\]
woraus sich \(C''',B''',A'''\) ergeben. Und so fort.

Beide Arten sind ungefähr gleich bequem, wenn die Gewichte der Bestimmungen sämtlicher Größen \(x,y,z\) etc. gesucht sind; werden dagegen nur eine oder zwei der Größen \([\alpha\alpha],\allowbreak[\beta\beta],\allowbreak[\gamma\gamma]\) etc. benötigt, so ist offenbar das erste System weit vorzuziehen.

Im Übrigen führt die Verbindung der Gleichungen I mit IV zu denselben Formeln und liefert außerdem zwei Rechnungen zur Ermittlung der plausibelsten Werte \(A,B,C\) etc. selbst, nämlich \textit{erstens}
\[
\begin{aligned}
A&=-\frac{\mathfrak{L}^0}{\mathfrak{A}^0}-A'\frac{\mathfrak{L}'}{\mathfrak{B}'}-A''\frac{\mathfrak{L}''}{\mathfrak{C}''}-A'''\frac{\mathfrak{L}'''}{\mathfrak{D}'''}-\text{etc.}\\[4pt]
B&=-\frac{\mathfrak{L}'}{\mathfrak{B}'}-B''\frac{\mathfrak{L}''}{\mathfrak{C}''}-B'''\frac{\mathfrak{L}'''}{\mathfrak{D}'''}-\text{etc.}\\[4pt]
C&=-\frac{\mathfrak{L}''}{\mathfrak{C}''}-C'''\frac{\mathfrak{L}'''}{\mathfrak{D}'''}-\text{etc.}\\
&\text{etc.}
\end{aligned}
\]
Die andere Rechnung ist mit der gewöhnlichen identisch, bei der \(u^0=0,\allowbreak u'=0,\allowbreak u''=0\) etc. gesetzt wird.

\Article{34}
Was wir in Art.\ 32 dargelegt haben, sind nur Spezialfälle eines allgemeineren Satzes, der wie folgt lautet:

\textsc{Satz.} \textit{Es bezeichne \(t\) die folgende lineare Funktion der Unbestimmten \(x,y,z\) etc.}
\[
t=fx+gy+hz+\text{etc.}+k
\]
\textit{die bei ihrer Umformung in eine Funktion der Unbestimmten \(u^0,u',u''\) etc. übergeht in}
\[
t=k^0u^0+k'u'+k''u''+\text{etc.}+K
\]
\textit{Dann ist \(K\) der plausibelste Wert von \(t\) und das Gewicht dieser Bestimmung ist}
\[
=\frac1{\mathfrak{A}^0k^0k^0+\mathfrak{B}'k'k'+\mathfrak{C}''k''k''+\text{etc.}}
\]
\textit{Beweis.} Der erste Teil des Satzes ist klar, weil der plausibelste Wert von \(t\) den Werten \(u^0=0,\allowbreak u'=0,\allowbreak u''=0\) etc. entsprechen muss. Zum Beweis des zwei\SourceLeaf{047}%
ten Teils bemerken wir: Wegen \(\tfrac12d\Omega=\xi\,dx+\eta\,dy+\zeta\,dz+\text{etc.}\) und \(dt=f\,dx+g\,dy+h\,dz+\text{etc.}\) gilt für \(\xi=f.\quad\eta=g,\allowbreak\zeta=h\) etc., unabhängig von den Werten der Differentiale \(dx,dy,dz\) etc.,
\[
d\Omega=2dt
\]
Daraus folgt für dieselben Werte \(\xi=f,\allowbreak\eta=g,\allowbreak\zeta=h\) etc.
\[
\begin{aligned}
&\frac{u^0}{\mathfrak{A}^0}du^0+\frac{u'}{\mathfrak{B}'}du'+\frac{u''}{\mathfrak{C}''}.du''+\text{etc.}\\
&\qquad=k^0du^0+k'du'+k''du''+\text{etc.}
\end{aligned}
\]
Nun sieht man leicht, dass, wenn \(dx,dy,dz\) etc. voneinander unabhängig sind, auch \(du^0,du',du''\) etc. voneinander unabhängig sind; daraus schließen wir, dass für \(\xi=f,\allowbreak\eta=g,\allowbreak\zeta=h\) etc. gilt
\[
u^0=\mathfrak{A}^0k^0,\quad u'=\mathfrak{B}'k',\quad u''=\mathfrak{C}''k''\quad\text{etc.}
\]
Der denselben Werten entsprechende Wert von \(\Omega\) ist daher
\[
=\mathfrak{A}^0k^0k^0+\mathfrak{B}'k'k'+\mathfrak{C}''k''k''+\text{etc.}+M
\]
woraus nach Art.\ 29 die Richtigkeit unseres Satzes unmittelbar folgt.

Wollen wir im Übrigen die Umformung der Funktion \(t\) unmittelbar, d.\ h. ohne Kenntnis der Einsetzungen IV aus Art.\ 32, ausführen, so stehen die Formeln
\[
\begin{aligned}
f&=\mathfrak{A}^0k^0\\
g&=\mathfrak{B}^0k^0+\mathfrak{B}'k'\\
h&=\mathfrak{C}^0k^0+\mathfrak{C}'k'+\mathfrak{C}''k''\quad\text{etc.,}
\end{aligned}
\]
zur Verfügung, aus denen die Koeffizienten \(k^0,k',k''\) etc. nacheinander bestimmt werden; schließlich haben wir
\[
K=k-\mathfrak{L}^0k^0-\mathfrak{L}'k'-\mathfrak{L}''k''-\text{etc.}
\]

\clearpage
